% GENERATED FILE. Edit canonical lessons, then run npm run generate:books.
% canonical-source-hash: fnv1a64:a858c99bac6f3605
% canonical-lessons: ZH-C428-tongwang, ZH-C428-qianwang, ZH-C428-huicheng, ZH-C428-xingcheng, ZH-C428-fapiao

\chapter{Moving On}
\label{ch:zh-moving-on}
\hlchaptermodality{428}

\begin{chapteropening}
\textbf{By the end of this chapter:} \emph{I can say where a door leads, follow an announcement, and talk about the trip back, the plan and receipts.}

Complete the last lesson of chapter 428: I can say where a door leads, follow an announcement, and talk about the trip back, the plan and receipts.
\end{chapteropening}

\section[\texorpdfstring{\zh{通往} (tōngwǎng)}{通往 (tōngwǎng)}]{\zh{通往} (tōngwǎng) — to lead to}

\label{lesson:ZH-C428-tongwang}



\begin{quote}
Before the new one: say the Mandarin for a bed sheet, then the Mandarin for to be bound for a place.
\end{quote}

\subsection*{You'll want to know: \zh{通往}}
\textbf{\zh{通往}} — \emph{tōngwǎng} — \textquotedblleft{}to lead to\textquotedblright{}.

\begin{quote}
\textbf{\zh{这个门通往停车场}} \emph{zhège mén tōngwǎng tíngchēchǎng} — This door leads to the car park.
\end{quote}

\begin{cousinweb}
\textbf{\zh{通往}} (\emph{tōngwǎng}) is built from \zh{通} 'to go through' and \zh{往} 'to go towards', both inherited from Middle Chinese, \emph{thuwng} and \emph{hjwangX} (Guangyun, in Baxter's transcription).
\end{cousinweb}

\begin{grammarlens}[title={What you've built}]
One more word for getting around.
\end{grammarlens}

\subsection*{Your turn}
\begin{itemize}
\raggedright
  \item \emph{Say it:} \emph{tōngwǎng}
  \item \emph{Say it:} \emph{tōngwǎng}, once more
  \item \emph{Say it:} \emph{zhège mén tōngwǎng tíngchēchǎng} — This door leads to the car park.
  \item \emph{Recall:} say the Mandarin for a bed sheet, then the Mandarin for to be bound for a place, then say \emph{tōngwǎng} again
\end{itemize}

\begin{tcolorbox}[breakable,colback=teal!4,colframe=teal!35!black,title={Before you move on}]
What does \zh{通往} mean? (\textquotedblleft{}To lead to\textquotedblright{}.) Say it once more.
\end{tcolorbox}
\section[\texorpdfstring{\zh{前往} (qiánwǎng)}{前往 (qiánwǎng)}]{\zh{前往} (qiánwǎng) — to set out for a place}

\label{lesson:ZH-C428-qianwang}



\begin{quote}
Before the new one: say the Mandarin for to be bound for a place, then the Mandarin for to lead to a place.
\end{quote}

\subsection*{You'll want to know: \zh{前往}}
\textbf{\zh{前往}} — \emph{qiánwǎng} — \textquotedblleft{}to set out for a place\textquotedblright{}.

\begin{quote}
\textbf{\zh{前往机场的旅客请上车}} \emph{qiánwǎng jīchǎng de lǚkè qǐng shàng chē} — Passengers for the airport, please get on.
\end{quote}

\begin{cousinweb}
\textbf{\zh{前往}} (\emph{qiánwǎng}) is built from \zh{前} 'forward' and \zh{往} 'to go', both inherited from Middle Chinese, \emph{dzen} and \emph{hjwangX} (Guangyun, in Baxter's transcription), which glosses the first 'before'.
\end{cousinweb}

\begin{grammarlens}[title={What you've built}]
One more word for getting around.
\end{grammarlens}

\subsection*{Your turn}
\begin{itemize}
\raggedright
  \item \emph{Say it:} \emph{qiánwǎng}
  \item \emph{Say it:} \emph{qiánwǎng}, once more
  \item \emph{Say it:} \emph{qiánwǎng jīchǎng de lǚkè qǐng shàng chē} — Passengers for the airport, please get on.
  \item \emph{Recall:} say the Mandarin for to be bound for a place, then the Mandarin for to lead to a place, then say \emph{qiánwǎng} again
\end{itemize}

\begin{tcolorbox}[breakable,colback=teal!4,colframe=teal!35!black,title={Before you move on}]
What does \zh{前往} mean? (\textquotedblleft{}To set out for a place\textquotedblright{}.) Say it once more.
\end{tcolorbox}
\section[huíchéng]{huíchéng — return journey; the journey back}

\label{lesson:ZH-C428-huicheng}



\begin{quote}
Before the new one: say the Mandarin for to lead to a place, then the Mandarin for to set out for a place.
\end{quote}

\subsection*{You'll want to know: huíchéng}
\emph{huíchéng} — \textquotedblleft{}return journey; the journey back\textquotedblright{}.

\begin{quote}
\emph{huíchéng wǒmen zuò gāotiě} — On the way back we're taking the fast train.
\end{quote}

Its characters come later; for now it is a word for the ear.

\begin{cousinweb}
\emph{huíchéng} is built from \zh{回} 'to go back' and \emph{chéng} 'a stage of a journey', both inherited from Middle Chinese, \emph{hwoj} and \emph{drjeng} (Guangyun, in Baxter's transcription), which glosses the second 'a set period' and 'a rule'; the sense 'a stretch of road' is not given there.
\end{cousinweb}

\begin{grammarlens}[title={What you've built}]
One more word for getting around.
\end{grammarlens}

\subsection*{Your turn}
\begin{itemize}
\raggedright
  \item \emph{Say it:} \emph{huíchéng}
  \item \emph{Say it:} \emph{huíchéng}, once more
  \item \emph{Say it:} \emph{huíchéng wǒmen zuò gāotiě} — On the way back we're taking the fast train.
  \item \emph{Recall:} say the Mandarin for to lead to a place, then the Mandarin for to set out for a place, then say \emph{huíchéng} again
\end{itemize}

\begin{tcolorbox}[breakable,colback=teal!4,colframe=teal!35!black,title={Before you move on}]
What does huíchéng mean? (\textquotedblleft{}Return journey; the journey back\textquotedblright{}.) Say it once more.
\end{tcolorbox}
\section[xíngchéng]{xíngchéng — itinerary; plan for a trip}

\label{lesson:ZH-C428-xingcheng}



\begin{quote}
Before the new one: say the Mandarin for to set out for a place, then the Mandarin for the journey back.
\end{quote}

\subsection*{You'll want to know: xíngchéng}
\emph{xíngchéng} — \textquotedblleft{}itinerary; plan for a trip\textquotedblright{}.

\begin{quote}
\emph{wǒ zhǎo bu dào wǒmen de xíngchéng le} — I can't find our itinerary.
\end{quote}

Its characters come later; for now it is a word for the ear.

\begin{cousinweb}
\emph{xíngchéng} is built from \emph{xíng} 'to go' and \emph{chéng} 'a stage', both inherited from Middle Chinese, \emph{haeng} and \emph{drjeng} (Guangyun, in Baxter's transcription).
\end{cousinweb}

\begin{grammarlens}[title={What you've built}]
One more word for getting around.
\end{grammarlens}

\subsection*{Your turn}
\begin{itemize}
\raggedright
  \item \emph{Say it:} \emph{xíngchéng}
  \item \emph{Say it:} \emph{xíngchéng}, once more
  \item \emph{Say it:} \emph{wǒ zhǎo bu dào wǒmen de xíngchéng le} — I can't find our itinerary.
  \item \emph{Recall:} say the Mandarin for to set out for a place, then the Mandarin for the journey back, then say \emph{xíngchéng} again
\end{itemize}

\begin{tcolorbox}[breakable,colback=teal!4,colframe=teal!35!black,title={Before you move on}]
What does xíngchéng mean? (\textquotedblleft{}Itinerary; plan for a trip\textquotedblright{}.) Say it once more.
\end{tcolorbox}
\section[\texorpdfstring{\zh{发票} (fāpiào)}{发票 (fāpiào)}]{\zh{发票} (fāpiào) — official receipt}

\label{lesson:ZH-C428-fapiao}



\begin{quote}
Before the new one: say the Mandarin for the journey back, then the Mandarin for the plan for a trip.
\end{quote}

\subsection*{You'll want to know: \zh{发票}}
\textbf{\zh{发票}} — \emph{fāpiào} — \textquotedblleft{}official receipt\textquotedblright{}.

\begin{quote}
\textbf{\zh{我的发票不见了}} \emph{wǒ de fāpiào bú jiàn le} — My receipt has gone missing.
\end{quote}

\begin{cousinweb}
\textbf{\zh{发票}} (\emph{fāpiào}) is built from \emph{fā} 'to send out' and \zh{票} 'a slip'. The first is inherited from Middle Chinese \emph{pjot} (Guangyun, in Baxter's transcription), glossed 'to start' and 'to send out'; no dictionary we use records an older form of \zh{票}, so its history is not settled.
\end{cousinweb}

\begin{grammarlens}[title={What you've built}]
One more word for getting around.
\end{grammarlens}

\subsection*{Your turn}
\begin{itemize}
\raggedright
  \item \emph{Say it:} \emph{fāpiào}
  \item \emph{Say it:} \emph{fāpiào}, once more
  \item \emph{Say it:} \emph{wǒ de fāpiào bú jiàn le} — My receipt has gone missing.
  \item \emph{Recall:} say the Mandarin for the journey back, then the Mandarin for the plan for a trip, then say \emph{fāpiào} again
\end{itemize}

\begin{tcolorbox}[breakable,colback=teal!4,colframe=teal!35!black,title={Before you move on}]
What does \zh{发票} mean? (\textquotedblleft{}Official receipt\textquotedblright{}.) Say it once more.
\end{tcolorbox}
